% ID: SEL-D0027-OA-E02
% ESTADO: validada. S028, cotejada con examen y criterios.
\textbf{a)} $f'(x)=3x^2-12=3(x-2)(x+2)$.
La derivada es positiva fuera de $[-2,2]$ y negativa en $(-2,2)$.
Por tanto, $f$ crece hasta $-2$, decrece entre $-2$ y $2$,
y vuelve a crecer a partir de $2$. El cambio de signo de
$+$ a $-$ da un máximo relativo en $\boxed{(-2,17)}$.
El cambio inverso da un mínimo relativo en $\boxed{(2,-15)}$.

\textbf{b)} $f(1)=-10$ y $f'(1)=-9$. La tangente tiene
ecuación $y+10=-9(x-1)$, es decir, $\boxed{y=-9x-1}$.
